% Part: sets-functions-relations
% Chapter: functions
% Section: isomorphisms

\documentclass[../../../include/open-logic-section]{subfiles}

\begin{document}

\olfileid{sfr}{fun}{iso}
\olsection{సమరూపత}

\begin{explain}
\emph{సమరూపణం} అనేది తాను అనుసంధానించే సమితుల నిర్మాణాన్ని నిలుపుకునే ద్వైక్యం. ఇక్కడ నిర్మాణం అంటే ఆ సమితుల \tetoken{మూలకాల}{!!{element}s} మధ్య ఉండే సంబంధాలు. $X=\{1,2,3\}$, $Y=\{4,5,6\}$ అనే రెండు సమితులను చూడండి. ఈ రెండింటికీ ఉత్తరవర్తి, కంటే చిన్నది, కంటే పెద్దది అనే సంబంధాల ద్వారా నిర్మాణం ఏర్పడుతుంది. ఆ నిర్మాణాలను నిలుపుకునే ఒక \tetoken{ద్వైక్యం}{!!a{bijection}} ఈ సమితుల మధ్య సమరూపణం. అంటే ఒక \tetoken{ద్వైక్య}{!!a{bijective}} ప్రమేయం $f \colon X \to Y$ సమరూపణం అవడానికి కింది నిబంధనలు అవసరమైనవీ చాలినవీ: $i<j$ అవడం $f(i)<f(j)$ అవడానికి అవసరమైనదీ చాలినదీ; $i>j$ అవడం $f(i)>f(j)$ అవడానికి అవసరమైనదీ చాలినదీ; $j$ అనేది $i$ యొక్క ఉత్తరవర్తి అవడం $f(j)$ అనేది $f(i)$ యొక్క ఉత్తరవర్తి అవడానికి అవసరమైనదీ చాలినదీ.
\end{explain}

\begin{defn}[సమరూపణం]
$U$ అనేది $\langle X, R\rangle$ అనే జంట, $V$ అనేది $\langle Y, S\rangle$ అనే జంట అనుకుందాం. ఇక్కడ $X$, $Y$ సమితులు; $R$, $S$ వరుసగా $X$, $Y$లపై సంబంధాలు. $X$ నుంచి $Y$కు ఒక \tetoken{ద్వైక్యం}{!!{bijection}} $f$ సంబంధ నిర్మాణాన్ని నిలుపుకుంటే, నిలుపుకున్నప్పుడే, అది $U$ నుంచి $V$కు \emph{సమరూపణం} అవుతుంది. అంటే $X$లోని ఏ $x_{1}$, $x_{2}$లకైనా $\tuple{x_1,x_2} \in R$ అవడం $\tuple{f(x_1),f(x_2)} \in S$ అవడానికి అవసరమైనదీ చాలినదీ.
\end{defn}

\begin{ex}
$X=\{1,2,3\}$, $Y=\{4,5,6\}$ అనే రెండు సమితులను, కంటే చిన్నది, కంటే పెద్దది అనే సంబంధాలను చూడండి. $f(x) = 7-x$ ద్వారా ఇచ్చిన ప్రమేయం $f\colon X \to Y$ అనేది $\tuple{X,<}$, $\tuple{Y,>}$ల మధ్య సమరూపణం.
\end{ex}

\end{document}
